\documentclass[a4paper, serif, 10pt]{moderncv}

%% moderncv
% color and style
\moderncvcolor{black}
\moderncvstyle{banking}

%% page layout/margins
\usepackage[top=2.5cm, bottom=2.5cm, left=1.5cm, right=1.5cm]{geometry}

\nopagenumbers{}

%% Babel config for German language
% ref:
% - https://latex3.github.io/babel/guides/locale-german.html
\usepackage[ngerman]{babel}

%% fontspec
\usepackage{fontspec}

\IfFontExistsTF{Linux Libertine}{
  \setmainfont[Ligatures={TeX, Common}, Numbers=OldStyle]{Linux Libertine}
}{}
\IfFontExistsTF{Linux Libertine O}{
  \setmainfont[Ligatures={TeX, Common}, Numbers=OldStyle]{Linux Libertine O}
}{}

%% CTeX
% CJK support, used to show CJK characters in the resume
%
% - heading=false: disable ctex heading
% - punct=kaiming: use kaiming punctuations styles for CJK
% - scheme=plain: use plain scheme, do not override \normalsize font size
% - space=auto: space settings for CJK characters
%
% CJK fonts are resolved per locale language, e.g. Noto Serif/Sans CJK JP
% for Japanese resumes, with fallback to the next region if not installed
%
% ref:
% - http://ctan.mirrorcatalogs.com/language/chinese/ctex/ctex.pdf
\usepackage[UTF8, heading=false, punct=kaiming, scheme=plain, space=auto]{ctex}

\IfFontExistsTF{Noto Serif CJK SC}{
  \setCJKmainfont{Noto Serif CJK SC}
}{}
\IfFontExistsTF{Noto Sans CJK SC}{
  \setCJKsansfont{Noto Sans CJK SC}
}{}

%% line spacing
\usepackage{setspace}
\setstretch{1.125}

%% URL styling - use normal text instead of monospace
\urlstyle{same}

%% Patch \httplink and \httpslink to handle full URLs
% The original moderncv macros always prepend http:// or https:// to URLs.
% This causes issues when the URL already has a protocol (e.g., https://example.com)
% resulting in malformed URLs like https://https://example.com.
% This patch checks if the URL already contains "://" and uses it directly if so.
% Note: We use \str_set:Nx (x-expansion) to expand macros like \@homepage first.
\ExplSyntaxOn
\str_new:N \g_yamlresume_url_str

\RenewDocumentCommand{\httpslink}{O{}m}{
  \str_set:Nx \g_yamlresume_url_str {#2}
  \str_if_in:NnTF \g_yamlresume_url_str {://}
    { \url{#2} }
    { \url{https://#2} }
}

\RenewDocumentCommand{\httplink}{O{}m}{
  \str_set:Nx \g_yamlresume_url_str {#2}
  \str_if_in:NnTF \g_yamlresume_url_str {://}
    { \url{#2} }
    { \url{http://#2} }
}
\ExplSyntaxOff

\name{Lena Müller}{}
\title{QA Engineer}
\phone[mobile]{+49 30 12345678}
\email{lena.mueller@posteo.de}
\homepage{https://www.lena-mueller-qa.de}

\address{Musterstraße 12, Berlin, Berlin, Deutschland, 10115}{}{}

\extrainfo{{\small \faLinkedin}\ \href{https://www.linkedin.com/in/lena-mueller-qa}{@lena-mueller-qa} {} {} {} • {} {} {} 
{\small \faGithub}\ \href{https://github.com/lena-mueller}{@lena-mueller}}

\begin{document}

\maketitle

\section{Basisinformationen}

\cvline{}{QA Engineer mit über 8 Jahren Erfahrung in der Qualitätssicherung von Web- und Cloud-Anwendungen. Spezialisiert auf Testautomatisierung, Testmanagement und agile Entwicklungsmethoden.
%
\begin{itemize}
\item Entwicklung und Umsetzung von Teststrategien für komplexe Systeme
\item Automatisierung von Regressionstests mit Python, Selenium und Cypress
\item Enge Zusammenarbeit mit Entwicklungsteams zur Sicherstellung hoher Softwarequalität
\item Zertifiziert nach ISTQB Advanced Level und Certified Scrum Master
\end{itemize}}
\section{Ausbildung}

\cventry{Okt. 2014–Sept. 2016}
        {Master, Informatik, Punkte: 1.7}
        {Technische Universität Berlin}
        {\url{https://www.tu.berlin.de/}}
        {}
        {\begin{itemize}
\item Schwerpunkt auf Softwaretechnik und Qualitätssicherung
\item Masterarbeit über modellbasiertes Testen
\end{itemize}
\textbf{Kurse}: Softwarequalität, Testmanagement, Automatisiertes Testen, Programmierparadigmen, Datenbanksysteme}
\section{Berufserfahrung}

\cventry{Jan. 2020–Heute}
        {Senior QA Engineer}
        {SAP SE}
        {\url{https://www.sap.com/}}
        {}
        {\begin{itemize}
\item Leitung des QA-Teams für die SAP S/4HANA Cloud
\item Entwicklung und Implementierung von Teststrategien für komplexe Microservices
\item Automatisierung von Regressionstests mit Python und Selenium
\item Zusammenarbeit mit Entwicklern und Product Ownern in agilen Sprints
\end{itemize}
\textbf{Schlüsselwörter}: Testautomatisierung, Selenium, Python, Agile, Microservices}

\cventry{Juli 2016–Dez. 2019}
        {QA Engineer}
        {Zalando SE}
        {\url{https://www.zalando.de/}}
        {}
        {\begin{itemize}
\item Manuelles und automatisiertes Testen von Webanwendungen und APIs
\item Erstellung von Testplänen und Testfällen basierend auf Anforderungen
\item Durchführung von Performance- und Lasttests mit JMeter
\item Fehlerverfolgung und -management mit JIRA
\end{itemize}
\textbf{Schlüsselwörter}: Manuelles Testen, API-Testing, JMeter, JIRA, Agile}
\section{Sprachen}

\cvline{Deutsch}{Muttersprache oder zweisprachig \hfill \textbf{Schlüsselwörter}: Muttersprache}
\cvline{Englisch}{Verhandlungssichere Kenntnisse \hfill \textbf{Schlüsselwörter}: C1, Business Fluent}
\section{Fähigkeiten}

\cvline{Testautomatisierung}{Experte \hfill \textbf{Schlüsselwörter}: Selenium, Cypress, Playwright, Robot Framework}
\cvline{Programmiersprachen}{Erfahren \hfill \textbf{Schlüsselwörter}: Python, Java, JavaScript, SQL}
\cvline{Testmanagement}{Erfahren \hfill \textbf{Schlüsselwörter}: ISTQB, Teststrategie, Risikobasiertes Testen, Testplanung}
\cvline{Agile Methoden}{Experte \hfill \textbf{Schlüsselwörter}: Scrum, Kanban, SAFe, Continuous Integration}
\section{Auszeichnungen}

\cventry{Dez. 2021}
        {SAP SE}
        {QA Team of the Year}
        {}
        {}
        {Auszeichnung für herausragende Leistungen bei der Einführung einer automatisierten Testsuite, die die Testabdeckung um 40\% erhöhte.}
\section{Zertifikate}

\cventry{März 2017}
        {ISTQB}
        {ISTQB Certified Tester Foundation Level}
        {\url{https://www.istqb.org/certifications/certified-tester-foundation-level}}
        {}
        {}

\cventry{Juni 2019}
        {ISTQB}
        {ISTQB Advanced Level Test Manager}
        {\url{https://www.istqb.org/certifications/advanced-level-test-manager}}
        {}
        {}

\cventry{Nov. 2020}
        {Scrum Alliance}
        {Certified Scrum Master}
        {\url{https://www.scrumalliance.org/certifications}}
        {}
        {}
\section{Projekte}

\cventry{Jan. 2021–Heute}
        {Ein leichtgewichtiges Python-Framework für API-Tests}
        {Open-Source Testframework}
        {\url{https://github.com/lena-mueller/testframework}}
        {}
        {\begin{itemize}
\item Entwicklung eines Open-Source-Tools zur Vereinfachung von API-Tests
\item Integration mit pytest und Unterstützung von Datengetriebenen Tests
\item Über 500 Sterne auf GitHub
\end{itemize}
\textbf{Schlüsselwörter}: Python, API, Testing}
\section{Interessen}

\cvline{Technologie}{Open Source, Künstliche Intelligenz, Cloud Computing}
\cvline{Sport}{Laufen, Yoga, Wandern}
\section{Ehrenamtliche Tätigkeiten}

\cventry{März 2018–Heute}
        {Mentorin für Softwaretesting}
        {Women in Tech Berlin}
        {\url{https://www.women-in-tech-berlin.de/}}
        {}
        {\begin{itemize}
\item Mentoring von Frauen, die in den Bereich Softwaretesting einsteigen möchten
\item Organisation von Workshops zu Testautomatisierung und ISTQB-Zertifizierung
\item Unterstützung bei der Karriereentwicklung
\end{itemize}}
    
\end{document}